\begin{proof}[Proof of \cref{obj:thm:unrestricted-mixed-count-upper}]
\begin{enumerate}
\item By \cref{obj:thm:unrestricted-uniform-mixed-count-certificate}, choose a universal constant \(0<\overline C<\infty\). For every \(n,d,q,P\) satisfying the stated domain, model, and sampling conditions, this constant gives
\[
E_P\!\left[
 \bigl(\widehat\tau^{\mathrm{mix}}_{n,d,q}(\mathbf O_n)-\Phi(P)\bigr)^2
\right]
\le \mathfrak U_{n,d,q}
\le \overline C\,r_{n,d,q},
\]
where \(\Phi(P)\) is the signed cell functional of
\cref{obj:def:cell-functional}. The envelope here has precisely the branch formulas in the statement.
% lean: CausalSmith.Stat.MarRareqLogfrontier.unrestricted_uniform_mixed_count_certificate

\item The model hypothesis permits application of
\cref{obj:lem:unrestricted-cell-identification}, which supplies
\[
\tau(P)=\Phi(P).
\]
Substituting this identity into the preceding risk bound yields
\[
E_P\!\left[
 \bigl(\widehat\tau^{\mathrm{mix}}_{n,d,q}(\mathbf O_n)-\tau(P)\bigr)^2
\right]
\le \mathfrak U_{n,d,q}
\le \overline C\,r_{n,d,q}.
\]
% lean: CausalSmith.Stat.MarRareqLogfrontier.unrestricted_cell_identification

\item To establish the auxiliary-risk comparison, fix an observed sample
\[
\mathbf o=(o_i)_{i=1}^n
\in
\bigl(\{0,\ldots,d-1\}\times\{0,1\}^{4}\bigr)^n,
\qquad
o_i=(X_i,A_i,S_i,R_i,R_iY_i).
\]
Define the prefix probabilities and remaining probability by
\[
w_{k_0}:=P_\omega(K_0=k_0)
=e^{-n/2}\frac{(n/2)^{k_0}}{k_0!}
\quad(k_0\in\mathbb Z_{\ge0}),
\qquad
\delta_{\mathrm{tail}}
:=1-\sum_{k_0=0}^{n}w_{k_0}
=P_\omega(K_0>n).
\]
The disjoint prefix events imply
\[
w_{k_0}\ge0,\qquad
\sum_{k_0=0}^{n}w_{k_0}\le1,\qquad
\delta_{\mathrm{tail}}\ge0.
\]
For each \(0\le k_0\le n\), define the assignment set and its output average by
\[
\Sigma_{k_0}
:=\bigl\{\sigma:\{1,\ldots,k_0\}\to\{1,2,3\}\bigr\},
\qquad
\overline\tau_{k_0}(\mathbf o)
:=3^{-k_0}\sum_{\sigma\in\Sigma_{k_0}}
 \widetilde\tau(\mathbf o;(k_0,\sigma)),
\]
where the three assignment values denote the membership, pilot, and estimation streams. The notation \((k_0,\sigma)\) specifies the auxiliary prefix and assignment in \cref{obj:def:mixed-count-estimator}.

The weights \(w_0,\ldots,w_n,\delta_{\mathrm{tail}}\) sum to one, and
\[
\sum_{k_0=0}^{n}w_{k_0}\overline\tau_{k_0}(\mathbf o)-\tau(P)
=
\delta_{\mathrm{tail}}\bigl(-\tau(P)\bigr)
+\sum_{k_0=0}^{n}w_{k_0}
 \bigl(\overline\tau_{k_0}(\mathbf o)-\tau(P)\bigr).
\]
Convexity of the square therefore gives
\[
\left(
 \sum_{k_0=0}^{n}w_{k_0}\overline\tau_{k_0}(\mathbf o)-\tau(P)
\right)^2
\le
\delta_{\mathrm{tail}}\tau(P)^2
+\sum_{k_0=0}^{n}w_{k_0}
 \bigl(\overline\tau_{k_0}(\mathbf o)-\tau(P)\bigr)^2.
\]
% lean: CausalSmith.Stat.MarRareqLogfrontier.finite_weighted_square_with_remainder

\item Since \(|\Sigma_{k_0}|=3^{k_0}\), centering the assignment average gives
\[
\overline\tau_{k_0}(\mathbf o)-\tau(P)
=
3^{-k_0}\sum_{\sigma\in\Sigma_{k_0}}
 \bigl(\widetilde\tau(\mathbf o;(k_0,\sigma))-\tau(P)\bigr).
\]
The finite mean-square inequality, obtained by Cauchy--Schwarz, consequently yields
\[
\bigl(\overline\tau_{k_0}(\mathbf o)-\tau(P)\bigr)^2
\le
3^{-k_0}\sum_{\sigma\in\Sigma_{k_0}}
 \bigl(\widetilde\tau(\mathbf o;(k_0,\sigma))-\tau(P)\bigr)^2.
\]
This also covers \(k_0=0\): the empty prefix has one assignment, so the averaging denominator equals one.
% lean: CausalSmith.Stat.MarRareqLogfrontier.finite_uniform_square

\item Multiplying these inequalities by the nonnegative prefix probabilities and summing gives
\[
\begin{aligned}
&\sum_{k_0=0}^{n}w_{k_0}
 \bigl(\overline\tau_{k_0}(\mathbf o)-\tau(P)\bigr)^2\\
&\qquad\le
\sum_{k_0=0}^{n}w_{k_0}3^{-k_0}
 \sum_{\sigma\in\Sigma_{k_0}}
 \bigl(\widetilde\tau(\mathbf o;(k_0,\sigma))-\tau(P)\bigr)^2.
\end{aligned}
\]
The estimator's finite averaging formula in
\cref{obj:def:mixed-count-estimator} is
\[
\widehat\tau^{\mathrm{mix}}_{n,d,q}(\mathbf o)
=
\sum_{k_0=0}^{n}w_{k_0}\overline\tau_{k_0}(\mathbf o).
\]
Combining this identity with the preceding two bounds establishes, for every \(\mathbf o\),
\[
\begin{aligned}
&\bigl(\widehat\tau^{\mathrm{mix}}_{n,d,q}(\mathbf o)-\tau(P)\bigr)^2\\
&\quad\le
\sum_{k_0=0}^{n}w_{k_0}3^{-k_0}
 \sum_{\sigma\in\Sigma_{k_0}}
 \bigl(\widetilde\tau(\mathbf o;(k_0,\sigma))-\tau(P)\bigr)^2
+\delta_{\mathrm{tail}}\tau(P)^2\\
&\quad=
E_\omega\!\left[
 \bigl(\widetilde\tau(\mathbf O_n;\omega)-\tau(P)\bigr)^2
 \,\middle|\,\mathbf O_n=\mathbf o
\right].
\end{aligned}
\]
The last term in the sum is the squared loss of the zero output on \(K_0>n\). All functions in this inequality are measurable and integrable on the finite observed-sample space. Integrating under the canonical product law yields
\[
E_P\!\left[
 \bigl(\widehat\tau^{\mathrm{mix}}_{n,d,q}(\mathbf O_n)-\tau(P)\bigr)^2
\right]
\le
E_P\!\left[
 E_\omega\!\left[
  \bigl(\widetilde\tau(\mathbf O_n;\omega)-\tau(P)\bigr)^2
  \,\middle|\,\mathbf O_n
 \right]
\right],
\]
as required.
% lean: CausalSmith.Stat.MarRareqLogfrontier.mixed_count_derandomization
\end{enumerate}
\end{proof}
